% ==============================================================
%  tracking_icons.tex
%  Tracking primitives for the Radar Night TikZ library.
%  Color role: geometryColor (purple) + resultsColor (orange)
% ==============================================================

% ---------------------------------------------------------------
% \Detection -- single radar detection with range circle
% ---------------------------------------------------------------
\newcommand{\Detection}{%
  \begin{scope}
    \draw[iconStroke, fill=white, rounded corners=2pt] (-0.55,-0.55) rectangle (0.55,0.55);
    \fill[radarGreen] (0,0) circle (0.045);
    \draw[radarGreen, line width=0.5pt, dashed] (0,0) circle (0.12);
  \end{scope}
}

% ---------------------------------------------------------------
% \Track -- smooth trajectory through detections
% ---------------------------------------------------------------
\newcommand{\Track}{%
  \begin{scope}
    \draw[iconStroke, fill=white, rounded corners=2pt] (-0.55,-0.55) rectangle (0.55,0.55);
    \draw[geometryColor, line width=1.2pt]
      (-0.35,-0.2) .. controls (-0.15,0.1) and (0.05,-0.1) .. (0.35,0.2);
    \foreach \p in {(-0.35,-0.2),(-0.2,0.0),(-0.05,0.08),(0.1,-0.05),(0.35,0.2)}{
      \fill[radarGreen] \p circle (0.04);
    }
  \end{scope}
}

% ---------------------------------------------------------------
% \GhostDetection -- single ghost target marker
% ---------------------------------------------------------------
\newcommand{\GhostDetection}{%
  \begin{scope}
    \draw[iconStroke, fill=white, rounded corners=2pt] (-0.55,-0.55) rectangle (0.55,0.55);
    \draw[radarOrange, line width=0.6pt] (0,0) circle (0.1);
    \draw[radarOrange, line width=0.4pt, dashed] (0,0) circle (0.16);
  \end{scope}
}

% ---------------------------------------------------------------
% \JPDAAssociation -- multiple tracks to multiple detections
% ---------------------------------------------------------------
\newcommand{\JPDAAssociation}{%
  \begin{scope}
    \draw[iconStroke, fill=white, rounded corners=2pt] (-0.55,-0.55) rectangle (0.55,0.55);
    % Tracks
    \foreach \x/\y/\c in {-0.35/0.2/1, -0.3/-0.15/2}{
      \fill[geometryColor] (\x,\y) circle (0.04);
    }
    % Detections
    \foreach \x/\y in {-0.1/0.15, -0.05/-0.1, 0.05/0.0, 0.2/0.1}{
      \fill[radarGreen] (\x,\y) circle (0.035);
    }
    % Association arrows (probabilistic)
    \draw[iconFaint, line width=0.3pt, ->] (-0.35,0.2) -- (-0.1,0.15);
    \draw[iconFaint, line width=0.3pt, ->] (-0.35,0.2) -- (-0.05,-0.1);
    \draw[iconFaint, line width=0.3pt, ->] (-0.3,-0.15) -- (-0.05,-0.1);
    \draw[iconFaint, line width=0.3pt, ->] (-0.3,-0.15) -- (0.05,0.0);
    \draw[iconFaint, line width=0.3pt, ->] (-0.3,-0.15) -- (0.2,0.1);
  \end{scope}
}

% ---------------------------------------------------------------
% \EKFPrediction -- EKF prediction step with covariance ellipse
% ---------------------------------------------------------------
\newcommand{\EKFPrediction}{%
  \begin{scope}
    \draw[iconStroke, fill=white, rounded corners=2pt] (-0.55,-0.55) rectangle (0.55,0.55);
    % Previous state
    \fill[iconFaint] (-0.3,-0.1) circle (0.035);
    % Prediction arrow
    \draw[geometryColor, -{Latex[length=1.5mm]}, line width=0.8pt] (-0.3,-0.1) -- (0.0,0.1);
    % Predicted state
    \fill[geometryColor] (0.0,0.1) circle (0.04);
    % Covariance ellipse
    \draw[resultsColor, line width=0.6pt, dashed, rotate=30] (0.0,0.1) ellipse (0.12 and 0.06);
    % Measurement
    \draw[radarOrange, line width=0.5pt] (0.2,-0.15) circle (0.04);
    \fill[radarOrange] (0.2,-0.15) circle (0.025);
  \end{scope}
}

% ---------------------------------------------------------------
% \Trajectory -- complete trajectory with timestamps
% ---------------------------------------------------------------
\newcommand{\Trajectory}{%
  \begin{scope}
    \draw[iconStroke, fill=white, rounded corners=2pt] (-0.55,-0.55) rectangle (0.55,0.55);
    \foreach \i in {0,1,2,3}{
      \pgfmathsetmacro{\x}{-0.35 + \i*0.2}
      \pgfmathsetmacro{\y}{-0.1 + sin(\i*60)*0.2}
      \pgfmathsetmacro{\op}{0.5+\i*0.1}
      \fill[geometryColor, opacity=\op] (\x,\y) circle (0.035);
    }
  \end{scope}
}

% ---------------------------------------------------------------
% \MultipathScene -- clean multipath illustration for intro
% ---------------------------------------------------------------
\newcommand{\MultipathScene}{%
  \begin{scope}
    \draw[iconStroke, fill=white, rounded corners=2pt] (-0.55,-0.55) rectangle (0.55,0.55);
    
    % ---- ENVIRONMENT ----
    
    % Walls
    \draw[iconFaint, line width=0.8pt, dashed] (-0.48,-0.48) -- (0.48,-0.48);
    \draw[iconFaint, line width=0.8pt, dashed] (-0.48,-0.48) -- (-0.48,0.48);
    \draw[iconFaint, line width=0.8pt, dashed] (0.48,-0.48) -- (0.48,0.48);
    
    % ---- REAL PERSON (target) ----
    \foreach \p in {(-0.05,0.32),(0.0,0.28),(0.05,0.24),(0.0,0.2),(0.04,0.16),
                    (-0.04,0.16),(0.0,0.12),(0.06,0.08),(-0.06,0.08),(0.0,0.04),
                    (0.05,-0.02),(-0.05,-0.02),(0.03,-0.08),(-0.03,-0.08)}{
      \fill[radarGreen] \p circle (0.025);
    }
    
    % ---- RADAR ----
    \draw[iconStroke, fill=iconFaint, rounded corners=1pt] (-0.48,-0.1) rectangle (-0.35,0.0);
    \foreach \r in {0.05,0.09,0.13}{
      \draw[radarProcessing, line width=0.6pt] 
        (-0.35,-\r*0.5-0.05) arc (210:330:\r);
    }
    
    % ---- GHOST TARGETS (from reflections) ----
    \foreach \p in {(0.25,-0.3),(-0.2,0.2),(0.4,0.0)}{
      \draw[radarOrange, line width=0.6pt] \p circle (0.035);
      \draw[radarOrange, line width=0.4pt, dashed] \p circle (0.065);
    }
    
    % ---- REFLECTION PATHS ----
    % Direct path (solid green)
    \draw[radarGreen, line width=0.5pt, ->, opacity=0.5] 
      (-0.38,-0.05) -- (0.0,0.1);
    
    % Wall reflection (dashed orange)
    \draw[radarOrange, line width=0.4pt, dashed, ->, opacity=0.4] 
      (-0.38,-0.05) -- (0.2,-0.48) -- (0.0,0.1);
    
    % Furniture reflection (dotted orange)
    \draw[radarOrange, line width=0.4pt, dotted, ->, opacity=0.4] 
      (-0.38,-0.05) -- (0.4,-0.3) -- (0.0,0.1);
    
  \end{scope}
}

